\section{Uniform pinching on real two-planes}\label{sec:pinching}

We now prove the curvature bounds in Proposition~\ref{metric:transfer}.
The identity of Section~\ref{sec:curvature} reduces the problem to
finite-dimensional algebra.  The limiting main expression is nonnegative,
but it can vanish on certain real planes.
We identify those planes and show that the error with three horizontal
indices vanishes to the additional order needed there.

Throughout this section, horizontal forms are extended by zero on the
vertical factor of \(\mathbb C^2\oplus\mathbb C\).
We retain \(I,V,S_P,a_P,b_P,m_P,\Xi\) from
Section~\ref{sec:curvature}.  Unless otherwise indicated, constants depend
only on the fixed bounds \(cI\le h\le CI\) and the component bounds for
the error tensors.

\subsection{The mixed real-plane algebra}

The leading expression below can vanish even when both projections
of the real plane are nonzero.  For an explicit algebraic model, take
$h=I$, $E=2I$, and
$u_1=(e_1,1)$, $u_2=(ie_1,-i)$, where $e_1=(1,0)\in\mathbb C^2$.
Direct substitution gives
\[
 a_h+3b_h^2+4a_V+m_E+6b_Eb_V=4+4+4-12=0,
 \qquad \Xi_{\mathrm{mix}}=0.
\]
The horizontal and vertical projections have opposite complex
orientations.  The next lemma proves nonnegativity for all planes in
the stated range of $E$ and shows that every zero has this mixed-block
vanishing.

\begin{lemma}\label{pinch:algebra}
Let \(h\) be a positive definite horizontal Hermitian form, let \(E\) be
another horizontal Hermitian form with \(E\ge h\), and let \(u_1,u_2\)
be real-linearly independent vectors in \(\mathbb C^2\oplus\mathbb C\).
Define
\[
 P_h(E)=a_h+3b_h^2+4a_V+m_E+6b_Eb_V.
\]
If \(a_ha_V=0\), then \(P_h(E)>0\).
If \(a_ha_V>0\), put \(\tau=|b_h|/\sqrt{a_h}\).
For every real \(\beta\) such that \(E\le(2-\beta)h\),
\begin{equation}\label{pinch:square}
 P_h(E)\ge
 \left(\sqrt{a_h+3b_h^2}-2\sqrt{a_V}\right)^2
 +2\beta(1+\tau)\sqrt{a_ha_V}.
\end{equation}
In particular, \(h\le E\le2h\) implies \(P_h(E)\ge0\).
If equality holds in this last conclusion, then
\[
 a_h=a_V>0,\qquad |b_h|=\sqrt{a_h},\qquad b_Eb_V<0,
\]
and both mixed blocks \(\Xi^{i\bar0}\) and \(\Xi^{0\bar i}\) vanish.
\end{lemma}

\begin{proof}
If \(a_V=0\), then \(b_V=0\), because the vertical factor has complex
dimension one.
If \(a_h=0\), the two horizontal projections are linearly dependent
over \(\mathbb R\), since \(h\) is positive definite.
Their pairing under every Hermitian form is therefore real, so
\(b_E=0\).
Thus \(a_ha_V=0\) implies \(b_Eb_V=0\), and
\[
 P_h(E)\ge a_h+4a_V+m_h\ge a_{h+V}>0.
\]
Here the last inequality uses the positive definiteness of \(h+V\)
and the real linear independence of the pair.

Suppose now that \(A=a_h>0\) and \(U=a_V>0\).
The Hermitian Gram determinant gives \(0\le\tau\le1\).
Make a real change of basis in the plane so that its horizontal
projections, denoted \(u_H,v_H\), are real-orthonormal for \(h\).
Both ratios \(|b_h|/\sqrt{a_h}\) and \(|b_E|/\sqrt{a_h}\) are unchanged
by this operation.
The squared \(h\)-norms of \(u_H+i v_H\) and \(u_H-i v_H\) are
\(2(1+\tau)\) and \(2(1-\tau)\), in some order.
The absolute difference of their squared \(E\)-norms is \(4|b_E|\)
in the new basis.  Comparing one norm from above and the other from
below gives
\begin{equation}\label{pinch:imaginary-bound}
 \frac{|b_E|}{\sqrt A}
 \le\frac{(2-\beta)(1+\tau)-(1-\tau)}2
 =\frac{1+3\tau}{2}-\frac{\beta(1+\tau)}2
 \le\sqrt{1+3\tau^2}-\frac{\beta(1+\tau)}2.
\end{equation}
For either order of the two norms, the displayed right side is an
upper bound, since \(2-\beta\ge1\).
The last inequality follows from
\[
 1+3\tau^2-\left(\frac{1+3\tau}{2}\right)^2
 =\frac34(1-\tau)^2;
\]
it is strict unless \(\tau=1\).
This argument does not require \(\beta\ge0\).

The mixed Gram inequality \eqref{curv:mixed-gram} and the Hermitian
Gram inequality yield
\[
 m_E\ge2\sqrt{a_EU}\ge2|b_E|\sqrt U.
\]
Since \(U=b_V^2\), it follows that
\[
\begin{aligned}
 P_h(E)
 &\ge A+3b_h^2+4U-4|b_E|\sqrt U\\
 &\ge
 \left(\sqrt{A+3b_h^2}-2\sqrt U\right)^2
 +2\beta(1+\tau)\sqrt{AU},
\end{aligned}
\]
which proves \eqref{pinch:square}.

It remains to examine equality when \(h\le E\le2h\).
The zero-area case already treated shows that \(A,U>0\).
Set \(s=\sqrt{A+3b_h^2}\) and \(B=|b_E|\).
The inequalities above, with \(\beta=0\), give
\[
 P_h(E)\ge(s-2\sqrt U)^2+4\sqrt U(s-B)\ge0.
\]
Equality forces \(s=2\sqrt U\), \(B=s\), and equality in the bounds
used to reach this expression.
The strictness in \eqref{pinch:imaginary-bound} gives \(\tau=1\).
Consequently
\[
 U=A,\qquad B=2\sqrt A,\qquad b_Eb_V<0,\qquad
 m_E=2B\sqrt U.
\]
Since \(a_E\ge B^2\) and \(m_E\ge2\sqrt{a_EU}\), these identities force
\(a_E=B^2\) and equality in the mixed Gram inequality.
The horizontal real \(E\)-Gram matrix and the vertical real Gram matrix
are therefore proportional.

In the real \(h\)-orthonormalized basis used above, \(\tau=1\) says
that the horizontal projections have the form
\[
 u_H,\quad \sigma i u_H,\qquad \sigma\in\{1,-1\}.
\]
Their real \(E\)-Gram matrix is a positive scalar multiple of the
identity.  Proportionality implies that the vertical projections have
equal positive norms and real inner product zero.
The sign \(b_Eb_V<0\) selects the opposite complex orientation, so these
projections have the form
\[
 y,\quad-\sigma i y.
\]
Direct substitution now gives
\[
 \Xi^{i\bar0}
 =u_H^i\overline{-\sigma i y}-(\sigma i u_H^i)\bar y=0.
\]
The other mixed block vanishes as well.
Under any real change of basis in the plane, \(\Xi\) is multiplied by
the determinant of that change.  The vanishing therefore holds in the
original basis.
\end{proof}

\subsection{Uniform negativity for large \texorpdfstring{\(\lambda\)}{lambda}}

The next statement allows every error tensor with the specified bounds
and component types.  This uniformity will also give the compactness
argument needed after \(\lambda\) has been fixed.

\begin{lemma}\label{pinch:positivity}
Fix \(0<c\le C<\infty\) and \(0\le C_{\mathrm e}<\infty\).
There is a finite \(\lambda_0=\lambda_0(c,C,C_{\mathrm e})>0\) with the following
property.
Let \(\lambda\ge\lambda_0\), \(x\ge0\), and let \(h\) be a horizontal
Hermitian form with \(cI\le h\le CI\).  Put
\[
 q=x(1+x),\qquad G=\lambda I+xh,\qquad
 D=(1+2x)h-qhG^{-1}h.
\]
Let \(T_4,T_3\) be tensors with the Kähler symmetries and reality
condition, supported respectively on components with four and with
three horizontal indices, and satisfying
\[
 \|T_4\|\le C_{\mathrm e}x,\qquad
 \|T_3\|\le C_{\mathrm e}\sqrt q.
\]
Define
\[
 R=-\lambda S_I-qS_h-S_V-(S_{D+V}-S_D-S_V)+T_4+T_3.
\]
Then \(\mathcal H_R<0\) on every real two-plane.
\end{lemma}

\begin{proof}
Increase the prospective threshold so that \(\lambda\ge\max(1,C)\).
Then \eqref{curv:D-bounds} gives
\[
 (1+x)h\le D\le(1+2x)h.
\]
If the lemma were false, there would be a sequence
\(\lambda\to\infty\) of permitted data and real two-planes for which
\(F=-\mathcal H_R\le0\).
We suppress the sequence index.
By \eqref{curv:negative-numerator},
\begin{equation}\label{pinch:leading-numerator}
 F=\lambda(a_I+3b_I^2)+q(a_h+3b_h^2)+4a_V
 +m_D+6b_Db_V-\mathcal H_{T_4}-\mathcal H_{T_3}.
\end{equation}
The contracted estimates \eqref{curv:contracted-errors} apply to the
allowed error tensors by their component types and symmetries.

\paragraph{The regime \(q/\lambda\to0\).}
Suppose first that a subsequence has this property.
Monotonicity and linearity of the mixed Gram term give
\(m_D\ge c(1+x)m_I\), while
\[
 |b_D|\le\sqrt{a_D}\le C'(1+x)\sqrt{a_I}.
\]
Discarding nonnegative terms in \eqref{pinch:leading-numerator}, we get
\[
 F\ge\lambda a_I+4a_V+c(1+x)m_I
 -C'\bigl((1+x)\sqrt{a_Ia_V}+x a_I+
             \sqrt q\,\sqrt{a_Im_I}\bigr).
\]
Absorb the first product into \(2a_V\), and the last into
\(\frac c2(1+x)m_I\).  Since \(q/(1+x)=x\), this gives
\[
 F\ge[\lambda-C''((1+x)^2+x)]a_I+
 2a_V+\frac c2(1+x)m_I.
\]
Here \((1+x)^2+x=O(1+q)=o(\lambda)\), because \(\lambda\to\infty\).
Every coefficient is eventually positive, whereas
\(a_I+a_V+m_I=a_{I+V}>0\).  This contradicts \(F\le0\).

\paragraph{The regime with \(\lambda/q\) bounded.}
If the preceding regime has no subsequence, we may pass to one on
which \(q/\lambda\) is bounded below by a positive constant.
Then \(\lambda/q\) is bounded and \(x\to\infty\).
The term \(qS_h\) has four horizontal slots, whereas the vertical
term \(S_V\) already has coefficient one.  To balance these two blocks,
multiply horizontal vector components by \(q^{1/4}\).
In the resulting plane choose a real-orthonormal pair for
\(\operatorname{Re}(I+V)\).
The inverse transformation and a real basis change give a basis of
the original plane; the latter change multiplies \(F\) by a positive
factor.
Thus its sign is unchanged.
The rescaled pairs belong to a fixed compact Stiefel manifold, and
their quantities satisfy \(a_{I+V}=1\).

In the rescaled variables, the leading part of
\eqref{pinch:leading-numerator} is
\begin{equation}\label{pinch:rescaled-leading}
 L_x=\frac{\lambda}{q}(a_I+3b_I^2)+P_h(E),
 \qquad E=\frac D{\sqrt q}.
\end{equation}
The transformed tensors, denoted \(\widetilde T_4,\widetilde T_3\),
satisfy
\begin{equation}\label{pinch:rescaled-errors}
 \|\widetilde T_4\|\le C'\frac{x}{q}=O(x^{-1}),
 \qquad
 \|\widetilde T_3\|\le C'q^{-1/4}=O(x^{-1/2}).
\end{equation}
Indeed, each horizontal slot supplies a factor \(q^{-1/4}\).

The lower bound on \(D\) gives \(E\ge h\).
For an eigenvalue \(d\in[c,C]\) of \(h\), the corresponding ratio is
\[
 \frac Eh=\frac{1+2x}{\sqrt q}
           -\frac{\sqrt q\,d}{\lambda+xd}.
\]
Since
\[
 \frac{1+2x}{\sqrt q}=\sqrt{4+\frac1q}
 \le2+\frac{C_1}{x^2},
 \qquad
 \frac{\sqrt q\,d}{\lambda+xd}
 \ge c_1\frac{x}{\lambda+x},
\]
we have the uniform bound
\begin{equation}\label{pinch:E-bounds}
 h\le E\le
 \left(2+\frac{C_1}{x^2}
          -c_1\frac{x}{\lambda+x}\right)h.
\end{equation}
The constants \(c_1>0,C_1<\infty\) depend only on \(c,C\).

Pass to a subsequence on which \(h,E,\lambda/q\), and the normalized
pair converge.
The limiting forms satisfy \(h\le E\le2h\).
The error contributions tend to zero by
\eqref{pinch:rescaled-errors}.
Lemma~\ref{pinch:algebra} shows that the two limiting summands in
\eqref{pinch:rescaled-leading} are nonnegative.
Since \(F\le0\), both must vanish.
In particular \(P_h(E)=0\) at the limit, so \(a_h,a_V>0\) there and
the mixed blocks of the limiting \(\Xi\) vanish.
The comparison
\[
 c^2a_I\le a_h\le C^2a_I
\]
follows from the corresponding real Gram-matrix bounds.
Hence \(a_I,a_h,a_V\) are bounded below by positive constants along
this subsequence.

We now improve the error estimate on these particular planes.
By the contraction formula \eqref{curv:xi-contraction}, a tensor with exactly three
horizontal indices can pair only one horizontal-horizontal block
with one mixed block.  Thus
\begin{equation}\label{pinch:vanishing-error}
 |\mathcal H_{\widetilde T_3}|
 \le C'q^{-1/4}\|\Xi_{HH}\|\,\|\Xi_{\mathrm{mix}}\|
 =o(x^{-1/2}).
\end{equation}
Here \(\Xi_{HH}\) denotes the block \((\Xi^{i\bar j})\), and
\(\Xi_{\mathrm{mix}}\) comprises the blocks
\((\Xi^{i\bar0})\) and \((\Xi^{0\bar i})\).
The first block is bounded because the normalized pairs form a
compact set, and the mixed blocks tend to zero by the equality case
just proved.
The explicit tensor factor \(q^{-1/4}\) is what makes this an
\(o(x^{-1/2})\) estimate for the evaluated contribution.
No rate of convergence of the planes is being assumed.

Apply \eqref{pinch:square} before taking the limit, with
\[
 \beta=c_1\frac{x}{\lambda+x}-\frac{C_1}{x^2}.
\]
This choice is permitted by \eqref{pinch:E-bounds}, including when
\(\beta<0\).
The positive lower bounds for \(a_I,a_h,a_V\), and boundedness of the
normalized pairs, imply
\begin{equation}\label{pinch:positive-gap}
 L_x\ge c_2\left(\frac{\lambda}{q}
                  +\frac{x}{\lambda+x}\right)-\frac{C_2}{x^2}
\end{equation}
along the subsequence, for some \(c_2>0,C_2<\infty\).
For \(x\ge1\), if \(\lambda\ge x^{3/2}\), then
\[
 \frac{\lambda}{q}\ge\frac{x^{3/2}}{2x^2}
 =\frac1{2\sqrt x}.
\]
If instead \(\lambda\le x^{3/2}\), then
\[
 \frac{x}{\lambda+x}\ge\frac1{\sqrt x+1}
 \ge\frac1{2\sqrt x}.
\]
Thus \eqref{pinch:positive-gap} is bounded below by a positive
constant times \(x^{-1/2}\) for large \(x\).
The evaluated \(\widetilde T_4\) contribution is
\(O(x^{-1})=o(x^{-1/2})\), and
\eqref{pinch:vanishing-error} gives the same little-o comparison for
\(\widetilde T_3\).
It follows that \(F>0\) eventually, a contradiction.

The two regimes exhaust the possible failure sequences.
This proves the asserted finite threshold \(\lambda_0\), uniformly
over all permitted matrices, error tensors, and real planes.
\end{proof}

\subsection{Both curvature bounds at a fixed parameter}

\begin{proof}[Proof of Proposition~\ref{metric:transfer}]
The geometric conclusions concerning the domain and completeness were
proved in Lemma~\ref{metric:complete}.
It remains to obtain both uniform real sectional-curvature bounds.
By Proposition~\ref{curv:tensor}, the errors of the actual metric have
the types and bounds in Lemma~\ref{pinch:positivity}, with
\(C_{\mathrm e}\) depending only on the chart bounds of
Proposition~\ref{metric:transfer}.
Choose any finite \(\lambda\ge\lambda_0\).
Then \(-\mathcal H_R>0\) for every point and every real two-plane.
We show that the sectional-curvature quotient is bounded away from
zero and from negative infinity.

First restrict \(x\) to a bounded interval \([0,X]\).
Normalize the pairs for \(\operatorname{Re}(I+V)\).
The data consisting of \(x\), the matrix \(h\), the component arrays
of \(T_4,T_3\), and the normalized pair range over a closed bounded
subset of a finite-dimensional space: impose
\(cI\le h\le CI\), the stated component types and symmetries, and
the closed norm bounds \(C_{\mathrm e}x,C_{\mathrm e}\sqrt q\).
This set is compact.
The formulas for \(G,D,R\) are continuous on it.
Since Lemma~\ref{pinch:positivity} applies to this entire set,
\(-\mathcal H_R\) has a strictly positive minimum and a finite maximum.
The metric \(G+V\) is uniformly positive definite and bounded there,
so \(a_{G+V}\) is bounded above and away from zero.
Formula~\eqref{curv:sectional} therefore gives both a strictly negative
upper curvature bound and a finite lower curvature bound on
\(0\le x\le X\).

To cover the end \(x\to\infty\), use the horizontal rescaling from
the preceding proof.
Here \(\lambda\) remains fixed: the mixed form \(D/\sqrt q\) will
converge to \(h\), rather than to an arbitrary form between \(h\)
and \(2h\) as allowed in the preceding contradiction argument.
The rescaled Hermitian metric is
\[
 \widetilde g_x=G/\sqrt q+V.
\]
For the fixed finite \(\lambda\),
\begin{equation}\label{pinch:metric-limit}
 \frac G{\sqrt q}-h
 =\frac{\lambda}{\sqrt q}I+
   \left(\frac{x}{\sqrt q}-1\right)h
 =O_\lambda(x^{-1})
\end{equation}
uniformly in the permitted data.
Because \(G\) commutes with \(h\), the formula for \(D\) also gives
\[
 D=xh+\lambda(1+x)hG^{-1}.
\]
Consequently
\begin{equation}\label{pinch:D-limit}
 \frac D{\sqrt q}-h
 =\left(\frac{x}{\sqrt q}-1\right)h+
   \frac{\lambda(1+x)}{\sqrt q}hG^{-1}
 =O_\lambda(x^{-1}).
\end{equation}
The rescaled tensor identity now implies
\begin{equation}\label{pinch:tensor-limit}
 \widetilde R_x=-S_{h+V}+O_\lambda(x^{-1/2})
\end{equation}
in component norm, uniformly over all allowed data.
Indeed, the term \(\lambda S_I\) is divided by \(q\), the term
\(qS_h\) becomes exactly \(S_h\), the mixed term converges by
\eqref{pinch:D-limit}, and the error tensors satisfy
\eqref{pinch:rescaled-errors}.

For the limiting metric \(h+V\) and tensor \(-S_{h+V}\),
\[
 K=-\frac12\left(1+
     3\frac{b_{h+V}^2}{a_{h+V}}\right)
 \in[-2,-1/2].
\]
For an \((I+V)\)-real-orthonormal pair,
\(a_{h+V}\ge\min(c,1)^2\), so these area denominators are uniformly
bounded away from zero.  The forms \(h+V\) also range over a compact
bounded family.  The metric and tensor estimates
\eqref{pinch:metric-limit}--\eqref{pinch:tensor-limit} therefore imply
uniform convergence of the corresponding sectional-curvature
quotients over all real planes.
For sufficiently large \(X_\lambda<\infty\), they give, for example,
\[
 -3\le K_g\le-1/4\qquad(x\ge X_\lambda).
\]
Combining this with the compact-range bounds at \(X=X_\lambda\)
produces finite constants \(0<A\le B\) for which
\[
 -B\le K_g(\sigma)\le-A<0
\]
for every point and every real two-plane.
The tensor identity at \(x=0\), given in
\eqref{curv:zero-section}, is computed in a smooth normalized
fiber frame; equivalently the same bounds extend to the zero section
by smoothness of the metric.
This completes Proposition~\ref{metric:transfer}.
\end{proof}
